A public street-lighting network has the structure of the provisioning problem
on ships: LED drivers of several product families, installed in populations of
known size, divided among maintenance units that each hold their own spares, and
exposed to events that put a whole street out at once. This section applies the
method to such a network; every number in it is measured, from public records.

\subsection{Data}
The Department for Infrastructure of Northern Ireland publishes, under the Open
Government Licence, every fault report on its street lighting since 2020 and an
inventory of its about 300\,000 lighting points with the lantern type of each~\cite{dfi_lighting}.
An operational area plays the part of a ship, a lantern family the part of a
product reference and a month the part of a mission. Eight LED families are
distinguished, seven named product lines and a residual group, the smallest
with about 2\,000 points on the streets used. Fault reports carry the street but not the lantern, so only
streets equipped with a single family are used, where the attribution is exact;
they hold about 166\,000 LED points. A report of a single light out is taken as
demand for a lantern spare. Reports closed as cable or supply faults, or as no
fault found, are excluded, and a report within seven days of the previous one on
the same street is counted as the same repair, which removes about a third of
the reports. Exposure is the installed count times the dusk-to-dawn hours of
the month at the network's latitude, and the calendar month enters as a fixed
effect for seasonal reporting.

The inventory is a snapshot, and older lanterns were still being replaced by
LED during the period. The models are therefore fitted on April 2023 to
March 2025 with the snapshot of April 2025, a total of \niTrainFail{} failures
over \niAreas{} areas and \niFamilies{} families. Two tests are made: a split
within that period, fitting on its first year (\niWithinTrain{} failures) and
predicting the second (\niWithinTest{} failures, \niWithinMonths{} area-months),
and a forward test on April to October 2026 with the November 2026 snapshot
(\niForwardTest{} failures, \niForwardMonths{} area-months). Each test is run
once with all areas in the fit and once leaving each area out in turn, as a
holder never seen before.

\subsection{The Allocation, Measured}
Fitted with an area term, an area--family interaction and an area--month term,
the network gives an area-wide standard deviation of \niSigArea{} (profile
interval \niCiArea, boundary test $p=\niPArea$), an interaction of
\niSigInter{} (\niCiInter, $p\niPInter$) and a month term of \niSigMonth{}
(\niCiMonth, $p=\niPMonth$). Family rate ratios range from \niFamLo{} to
\niFamHi{} relative to the reference family. On this network the heterogeneity
between holders is almost entirely of the interaction kind: an area is not
worse for every family, but particular families do markedly better or worse in
particular areas. Under the aggregate criterion this is the allocation for which
risk pools best, and under the joint criterion the one that is hardest to
cover (Section~\ref{sec:multi}). The shared month term, the analogue of a
ship-wide event, is small but detectable.

\subsection{Plans Tested Out of Sample}
Table~\ref{tab:ni} gives the joint cover of a 99\,\% plan over the families an
area holds. In the within-period test with known areas, the Poisson plan covers
\niPoissonWithinKnownCover\,\% of area-months and the empirical-Bayes plan
\niEbGammaWithinKnownCover\,\%. The hierarchical predictive plan covers
\niHierPredWithinKnownCover\,\% with \niHierPredWithinKnownUnits{} units, against
\niPoissonWithinKnownUnits{} for Poisson and \niNbPooledWithinKnownUnits{} for
the pooled negative binomial, which reaches \niNbPooledWithinKnownCover\,\%.
Knowing which families run high in which area is what lets the hierarchical
plan reach the target without the pooled model's uniform margin.

For an area left out of the fit the picture changes. The interaction of an
unseen area is unknown, and with a standard deviation of \niSigInter{} the
hierarchical plan has to cover it, at \niHierPredWithinNewUnits{} units. The
pooled negative binomial, which spreads the same heterogeneity over all cells,
covers \niNbPooledWithinNewCover\,\% with \niNbPooledWithinNewUnits{} units,
while Poisson and empirical Bayes fall to \niPoissonWithinNewCover{} and
\niEbGammaWithinNewCover\,\%. For a new holder the large interaction makes the
hierarchical plan safe but expensive, and its history, once available, is
worth most of that stock.

The forward test fell in a quieter year: failure rates in April to October 2026
were below those of the fitting period, and every plan covered every
area-month, Poisson with \niPoissonForwardKnownUnits{} units and the
hierarchical predictive plan with \niHierPredForwardKnownUnits. It confirms that
the plans are not exceeded a year ahead but cannot separate the methods; a
fleet should expect drifts of this kind and refit as records accumulate.

\begin{table}[!t]
\caption{Northern Ireland street lighting: joint cover (\%) of a 99\,\% plan
over an area's families, mean units per area-month and mean log score per
family-month
(\textsc{measured}; new areas are left out of the fit in turn)}
\label{tab:ni}
\centering
\footnotesize
\setlength{\tabcolsep}{2.6pt}
\begin{tabular}{@{}lccccccc@{}}
\toprule
& \multicolumn{5}{c}{within-period test} & \multicolumn{2}{c}{forward test} \\
\cmidrule(lr){2-6}\cmidrule(lr){7-8}
& \multicolumn{3}{c}{known areas} & \multicolumn{2}{c}{new area} &
\multicolumn{2}{c}{known areas} \\
\cmidrule(lr){2-4}\cmidrule(lr){5-6}\cmidrule(lr){7-8}
Plan & cover & units & log sc. & cover & units & cover & units \\
\midrule
\tabinput{tables/tab_ni.tex}
\bottomrule
\end{tabular}
\end{table}
